
\begin{table}[!htbp] \centering 
  \caption{Effects of AI to whether or not they submit any answer on each task} 
  \label{tab:task_finish2} 
\small 
\begin{tabular}{@{\extracolsep{-5pt}}lcc} 
\\[-1.8ex]\hline 
\hline \\[-1.8ex] 
 & \multicolumn{2}{c}{\textit{Dependent variable:}} \\ 
\cline{2-3} 
\\[-1.8ex] & Task 1 Submitted & Task 2 Submitted \\ 
\\[-1.8ex] & (1) & (2)\\ 
\hline \\[-1.8ex] 
 GenAI Treatment Assigned (Trt) & 0.026 & 0.043 \\ 
  & (0.026) & (0.030) \\ 
 \hline \\[-1.8ex] 
Mean Y in Control Group & 0.88 & 0.83 \\ 
P-value (Trt) & 0.315 & 0.151 \\ 
Observations & 573 & 573 \\ 
R$^{2}$ & 0.014 & 0.016 \\ 
\hline 
\hline \\[-1.8ex] 
\end{tabular}
\\
\begin{minipage}{\textwidth}
{\footnotesize \emph{Notes}: 
           This table analyzes the effect of the treatment on the consultants submitting any answer to each question.
            Text of problems can be found in Appendix Section~\ref{sec:tasks}.
           All regressions include controls for gender, location, native english status, and low tenure.
              Reported entries are coefficient estimates with two-sided 95\% standard errors computed using Huber--White (HC0) robust variance.
            Significance stars and p-values indicate two-sided tests using Huber–White (HC0) robust variance:
* p \textless 0.05, ** p \textless 0.01, *** p \textless 0.001.
          \starlanguage}
\end{minipage}
\end{table}
